% Authored. Numbers come from generated/macros.tex.

\section{Conclusion}
\label{sec:conclusion}

Portugal's general-government balance, decomposed across institutional subsectors
over \PanelYears\ years, has a composition that is asymmetric and remarkably stable.
The Central Government balance is negative in every one of the
\CentralNegativeYears\ observations. The Social Security balance is positive in
\SsfPositiveYears, with a longest unbroken run of \SsfLongestPositiveRun\ years.
Regional and Local Government is the only subsector that changes sign with any
regularity, and it is small in both directions. The aggregate tracks Central
Government closely inside each statistical regime, correlating at
\HistoricalTrackingCorrelation\ before the 1995 splice and \ModernTrackingCorrelation\
after it. The aggregate is positive
in \NumPositiveYears\ years, and in each of them the combined non-Social-Security
balance is negative --- from which, given the identity, the rest of that year's
composition follows.

Against \BenchmarkCountries\ European reporters, the composition is distinctive in one
specific respect and unremarkable in another. A permanently deficit-running central
tier is common, so the Central Government finding is weaker than it appears in
isolation. The Social Security surplus sits in the upper tail on both frequency and
size. And Portugal is the only reporter whose every aggregate-surplus year combines
that surplus with a negative non-Social-Security balance --- a composition present in
\SurplusOffsettingShare\% of European surplus years, though Portugal's own count is
only \PtSurplusYears.

On the Social Security side, most of the recent increase in contributions is
algebraically associated with growth in the aggregate employee wage bill rather than with
movement in the contributions-to-wage-bill ratio. Of the \BaseContributionsChange\ million
euro rise in \BaseYear,
\[
\BaseContributionsChange = \BaseFromWageBill + \BaseFromRatio + \BaseRatioInteraction,
\]
the three terms being the wage-bill component, the ratio component and their interaction.
The wage bill itself rose by \BaseWageBillChange\ million over the same year, and that is
a second and separate identity:
\[
\BaseWageBillChange = \BaseFromEmployment + \BaseFromAverageWage + \BaseWageInteraction,
\]
in employees, average wage and interaction. The two share a factor but not a total, and
the terms of the second are components of $\Delta W$ rather than parts of the
\BaseFromWageBill\ million wage-bill component of the first. Read together they say that
the movement was predominantly in the reference base rather than in the ratio, and that
within the base it came more from wages per employee than from employee numbers. That is
a mechanism in the accounting sense --- it names the quantity the contributions followed
--- and not a causal account of why that quantity moved.

Two results cut against readings that the aggregate series invites.

The first concerns Central Government. Its balance is negative throughout the panel,
but its primary balance is positive in \PrimaryPositiveYears\ of the
\PrimaryObservedYears\ years for which detailed accounts exist ---
\HistoricalPrimaryPositive\ of \HistoricalPrimaryYears\ before the 1995 splice and
\ModernPrimaryPositive\ of \ModernPrimaryYears\ after it, so the positive years are if
anything a historical phenomenon rather than a recent one. A persistently negative B.9
does not imply a persistently negative primary balance, and interest is the arithmetic
that separates them: it averaged \HistoricalInterestMean\% of GDP in the historical
regime against \ModernInterestMean\% in the modern one.

The second concerns Social Security. The ESA 2010 Social Security Funds sector and
the CFP budget-execution systems are different accounting objects whose balances
differ in every one of the \SsfBoundaryYears\ years they overlap, by between
\SsfBoundaryMin\ and \SsfBoundaryMax\ million euro. Quoting one where the other
belongs introduces an error of that order with no signal that it has occurred.

\subsection{What this paper does not establish}

The decomposition is an accounting identity. It reallocates a recorded quantity and
estimates nothing, and three inferences do not follow from it.

It does not license a counterfactual. That a positive Social Security balance exceeds
a negative non-Social-Security balance in a given year says those two published
numbers stand in a particular ratio. It says nothing about what the aggregate would
have been under a different institutional arrangement, because that counterfactual
would move contribution rates, entitlements and transfers together and the identity
models none of them.

It does not attribute responsibility. A subsector accounting arithmetically for most
of a movement has not been shown to have produced it.

It does not speak to sustainability. Neither the headline nor the primary balance
determines the debt path, which also depends on stock-flow adjustments that are the
larger term in several years of these data.

It does not explain why Portugal's composition differs. The benchmark establishes that
it does, in a stated respect, while holding no institutional feature constant.

Two further limits are properties of the data rather than of the method. Magnitudes
are regime-specific and cannot be pooled across the 1995 splice; the pooled aggregate
mean of \GgMeanPooled\% of GDP describes neither the \GgMeanHistorical\% of
\PanelStart--1994 nor the \GgMeanModern\% of 1995--\PanelEnd. And
\ProvisionalYears\ are provisional at source: they are the years discussed in most
detail here and the years most likely to be restated by a later vintage.

\subsection{Extensions}

Two extensions follow naturally.

The contribution base of Section~\ref{sec:contributionbase} stops at the aggregate wage
bill. Splitting it by sector, contract type or earnings band would say which part of the
base carried the growth, and the same identity would apply unchanged; the data required
are published but are not in the bundle assembled here.

Extending the benchmark of Section~\ref{sec:benchmark} would sharpen it. It currently
compares compositions while holding nothing else constant, so it cannot separate a
Portuguese peculiarity from the consequence of an institutional choice made elsewhere:
whether a country operates a state-government tier, how it assigns contributory schemes
between tiers, how mature its pension system is, and how it routes transfers between
subsectors. Conditioning on those would turn a distribution into a comparison.
